\clearpage
\VerdierPhysicalPageAnchor{0185}%
\begingroup\small\setlength{\parskip}{0pt}\setlength{\emergencystretch}{1.5em}
\setlength{\abovedisplayskip}{4pt plus 2pt minus 2pt}%
\setlength{\belowdisplayskip}{4pt plus 2pt minus 2pt}%

\VerdierSemanticAnchor{ch3d-1-3-6-conditions-cont}
\noindent\textit{i) $F$ は $K(\mathcal A)$ を介して分解する。}

\noindent\textit{ii) $F$ はホモトピー同値を $C$ の同型に送る。}

\VerdierSemanticAnchor{ch3d-1-3-6-proof}
i) $\Rightarrow$ ii) は明らかである。ii) $\Rightarrow$ i) を示すには、
$F$ がホモトピックな二つの射を同じ射に送ることを示せば十分である。
$X$ を $\operatorname{comp}(\mathcal A)$ の対象とする。
$\operatorname{comp}(\mathcal A)$ 上の関手
\[
Y\longmapsto
\{\text{$X$ から $Y$ への射の対で、それらを結ぶホモトピーを備えたもの}\}
\]
は複体 $Z$ によって表現される：

\VerdierSemanticAnchor{ch3d-1-3-6-representing-complex}
\begin{equation}
\begin{aligned}
Z^i&=X^i\oplus X^i\oplus X^{i+1},\\[2pt]
d_Z^i&=
\begin{pmatrix}
d_X^i&0&\operatorname{id}_{X^{i+1}}\\
0&d_X^i&-\operatorname{id}_{X^{i+1}}\\
0&0&-d_X^{i+1}
\end{pmatrix}.
\end{aligned}
\tag{1.3.6.1}
\end{equation}
したがって、二つのホモトピックな射
$p_1,p_2\colon X\longrightarrow Z$ が存在する：

\VerdierSemanticAnchor{ch3d-1-3-6-morphisms}
\begin{equation}
p_1^i=
\begin{pmatrix}
\operatorname{id}_{X^i}\\0\\0
\end{pmatrix},
\qquad
p_2^i=
\begin{pmatrix}
0\\\operatorname{id}_{X^i}\\0
\end{pmatrix},
\tag{1.3.6.2}
\end{equation}
そして $F(p_1)=F(p_2)$ を示せば十分である。一方、恒等射の対
$(\operatorname{id}_X,\operatorname{id}_X)$ に自明なホモトピーを備えたものは、
射 $\pi\colon Z\longrightarrow X$ を定める：

\VerdierSemanticAnchor{ch3d-1-3-6-retraction}
\begin{equation}
\pi^i=(\operatorname{id}_{X^i},\operatorname{id}_{X^i},0)
\tag{1.3.6.3}
\end{equation}
そして $\pi p_1=\pi p_2=\operatorname{id}_X$ である。さらに、射 $\pi$、$p_1$、
$p_2$ がホモトピー同値であることは直ちに確かめられる。したがって
$F(p_1)$、$F(p_2)$、$F(\pi)$ は同型であり、
$F(\pi)F(p_1)=F(\pi)F(p_2)=\operatorname{id}_{F(X)}$ である。
ゆえに $F(p_1)=F(p_2)$ となり、補題の証明が完了する。

\VerdierSemanticAnchor{ch3d-1-3-7}
\medskip\noindent\textbf{注意 1.3.7.}\quad
命題 (1.3.5) により、$\mathcal A$ の導来圏は
$\operatorname{comp}(\mathcal A)$ の擬同型を可逆にするという普遍問題の解である
（第 II 章，2.2.3，b)）。
\endgroup
